\subsection{用語の定義}
本研究で用いる主要な用語を定義する．

\begin{defi}
ここに定義を記述する．
\end{defi}

\subsection{基礎理論}
研究内容を理解するために必要な基礎理論を記述する．数式は以下のように記述できる．

\begin{equation}
    y = f(x)
    \label{eq:sample}
\end{equation}

\subsection{図表の例}
図を挿入する場合は，画像ファイルを \texttt{images/} に配置し，以下のように記述する．

% \begin{figure}[H]
%     \centering
%     \includegraphics[width=0.8\textwidth]{images/sample.png}
%     \caption{図のキャプション}
%     \label{fig:sample}
% \end{figure}

表は以下のように記述する．

\begin{table}[H]
    \centering
    \caption{表のキャプション}
    \label{tab:sample}
    \begin{tabular}{|c|c|} \hline
        項目 & 値 \\ \hline
        A & 1 \\ \hline
        B & 2 \\ \hline
    \end{tabular}
\end{table}

